\chapter{Discussione, limiti e minacce alla validità}
\label{sec:discussion}

% Source notes: Project/docs/thesis_notes.md, checkpoint_2026-10-10.md (takeaways and caveats).

Questo capitolo interpreta i risultati alla luce delle cinque domande di ricerca, per poi indicare ciò che gli
esperimenti non sono in grado di dimostrare e quali scelte avrebbero potuto distorcerli.

\section{Quando conviene un System One specialista}
\label{sec:disc-synthesis}

\paragraph{Specializzazione (RQ1).} È l'addestramento sulla decisione a rendere GroundOne accurato. I modelli non
addestrati si collocano vicino al caso, mentre gli stessi modelli dopo l'addestramento raggiungono 0.769 e 0.832 di
accuratezza bilanciata a tre classi sull'insieme di test di RAGTruth, superando gli encoder specializzati in AUROC e un
System One generalista otto volte più grande nell'accuratezza a tre classi. Il vantaggio si concentra là dove si trovano i
dati di addestramento: su LLM-AggreFact il 3B è alla pari con il miglior encoder di piccole dimensioni, anziché superarlo,
e resta al di sotto del System One chiuso Jev e di un modello aperto da 27B secondo misure indipendenti; il 0.5B rimane
sotto l'encoder, e su VitaminC il 3B addestrato non fa meglio di quello non addestrato. Il confronto con un modello
decisionale generalista della stessa dimensione rende esplicito il compromesso. d1-3B, già addestrato a decidere ma su
un'ampia miscela di compiti, è nettamente dietro GroundOne su RAGTruth, soprattutto sulle risposte lunghe, alla pari con
il 3B su LLM-AggreFact e nettamente avanti su VitaminC. Uno specialista è quindi una scelta solida per il dominio su cui è
stato addestrato e un debole verificatore di fatti generale: la specializzazione garantisce accuratezza e robustezza sulle
risposte lunghe nel proprio dominio al prezzo del trasferimento, mentre ciò che migliora il trasferimento al di fuori di
esso è la dimensione del modello, oppure un'ampia miscela di addestramento.

\paragraph{Readout contro generazione (RQ2).} Il meccanismo System One contribuisce in due modi distinti. Il primo è il
costo: un solo passaggio sulla risposta sostituisce la generazione di $K$ verdetti, per cui il tempo risparmiato cresce
con la lunghezza della risposta, fino a un fattore di circa due rispetto a un giudice sintetico e di sei rispetto a uno
verboso per risposte di quattordici o più frasi. Il secondo, che questa tesi non si attendeva, è l'accuratezza. A parità
di pesi e di costo, decidere tra le tre lettere ammissibili è più accurato che prendere il token che il modello
genererebbe, di 0.25 per il 0.5B e di 0.07 per il 3B. Un System One addestrato va dunque inteso più come un classificatore
a tre classi posto sopra un modello linguistico che come un generatore che si arresta dopo un token, e questa distinzione
conta soprattutto per i modelli piccoli, che sono quelli alla portata di un portatile. Il confronto con i giudici
generativi, infine, misura anche l'addestramento: i giudici non sono addestrati, perciò le ampie differenze di accuratezza
tra essi e GroundOne vanno lette come l'effetto combinato della specializzazione e del readout (la lettura dei logit).

\paragraph{Calibrazione (RQ3).} Una temperatura e un bias stimati sui dati di sviluppo fanno sì che le probabilità
corrispondano alle frequenze osservate sull'insieme di test dello stesso dominio, e le probabilità calibrate rendono
efficace la predizione selettiva. Due lezioni vanno oltre GroundOne. In primo luogo, l'errore di calibrazione atteso, da
solo, non basta a ordinare i sistemi: dopo la stessa calibrazione, giudici le cui probabilità separano a malapena le
classi raggiungono lo stesso ECE di GroundOne, e soltanto l'AUROC e la curva rischio--copertura rivelano la differenza. In
secondo luogo, decisioni e probabilità devono essere ricavate in modo coerente: la classe calibrata più probabile
collassa verso la classe maggioritaria, mentre una soglia stimata sui dati di sviluppo recupera la decisione bilanciata e
mantiene l'etichetta coerente con la probabilità mostrata all'utente. La calibrazione è valida soltanto per la
distribuzione su cui è stata stimata, come mostrano i risultati su VitaminC.

\paragraph{Costo sul dispositivo (RQ4).} Il verificatore trova agevolmente posto nel portatile da 16\,GB:
l'addestramento raggiunge picchi di 2.6 e 5.9\,GB, e la verifica di una risposta lunga richiede 0.41 secondi con il 0.5B e
2.57 con il 3B, con un'energia di 5.5 e 33 joule. Il 0.5B rispetta un budget interattivo; il 3B, per una risposta tipica,
è più lento di un obiettivo interattivo e scambia tempo con accuratezza. Le misure di energia precisano i risultati sulla
velocità: un singolo prefill lungo assorbe più potenza della decodifica token per token, sicché l'energia risparmiata è
inferiore al tempo risparmiato, e la velocità non va interpretata come un risparmio equivalente di batteria. La memoria
necessaria quando il verificatore gira accanto a un generatore e a un retriever non è stata misurata.

\paragraph{Oltre un punteggio binario (RQ5).} L'output a tre classi separa bene le frasi contraddette da quelle non
menzionate nel data-to-text, dove RAGTruth concentra le proprie contraddizioni, e molto meno nel question answering e nella
sintesi, dove le contraddizioni sono rare nei dati di addestramento. Il 3B migliora considerevolmente questo aspetto, il
che suggerisce che al modello piccolo manchino contraddizioni più varie in addestramento, più che un'architettura diversa.
Poiché le politiche rigorosa e permissiva e la soglia derivano dalle stesse tre probabilità, un utente può modificare la
nozione di errore senza riaddestrare il modello, cosa che nessun verificatore binario consente.

\section{Limiti}
\label{sec:disc-limitations}

\paragraph{Una sola piattaforma e una sola famiglia di modelli.} Ogni risultato proviene da un unico portatile, un unico
runtime e un'unica famiglia di modelli. Latenza ed energia dipendono dall'hardware e dall'implementazione, e ciò che
dovrebbe trasferirsi sono i rapporti tra i metodi, non i valori assoluti. Non si effettua alcun confronto di latenza con
gli encoder o con il System One generalista, perché questi sono stati eseguiti in runtime diversi e in condizioni che
renderebbero il confronto iniquo.

\paragraph{Stabilità dell'addestramento.} Le due esecuzioni complete del 0.5B concordano su tutte le metriche aggregate,
ma una terza esecuzione è andata precocemente in collasso ed è stata interrotta. La procedura non prevede clipping del
gradiente e usa Adam senza correzione del bias, entrambe cause note di fine-tuning
instabile~\cite{pascanu2013difficulty,mosbach2021stability}, per cui ciascuna delle due potrebbe spiegare il collasso; nessuna delle
due ipotesi è stata verificata. Il 3B è stato addestrato una sola volta, su metà delle risposte di addestramento, e la sua
varianza rispetto al seed è ignota.

\paragraph{Dati di addestramento.} RAGTruth è l'unica fonte di addestramento. Le sue contraddizioni sono concentrate nel
data-to-text, e al di fuori di questo compito vengono apprese solo in parte. I modelli si affidano inoltre in parte alla
sovrapposizione lessicale: un supporto parafrasato appare non supportato, e una piccola modifica in una frase che per il
resto copia la fonte passa spesso inosservata. Questo punto cieco si riduce con la dimensione del modello, ma resta il
limite principale di entrambi i modelli.

\paragraph{Ambito della decisione.} La verifica avviene frase per frase e soltanto rispetto al documento recuperato. Una
frase vera nel mondo ma assente dal documento conta come non menzionata, il che è il comportamento corretto per la
groundedness ma non per la verifica dei fatti. Le frasi prive di contenuto verificabile sono escluse da una regola, e la
segmentazione stessa può dividere o unire le affermazioni in modo diverso da un lettore umano. Il verdetto su una frase
dipende soltanto dal documento e dalle frasi che la precedono.

\paragraph{Ambito statistico.} Gli intervalli bootstrap descrivono la variabilità dei dati di test, non quella
dell'addestramento, e circa cinquanta confronti sono stati sottoposti a test senza correzione per i test multipli. Le
conclusioni poggiano pertanto su effetti ben lontani dallo zero, e non si rivendicano differenze su LLM-AggreFact e
VitaminC inferiori alla dispersione tra seed. Diverse baseline sono state valutate su sottocampioni, il che ne allarga gli
intervalli.

\section{Minacce alla validità}
\label{sec:disc-threats}

\paragraph{Etichette.} La distinzione tra contraddetta e non menzionata si fonda sugli annotatori di RAGTruth, il cui
accordo sul tipo di allucinazione non è pubblicato, e una recente riannotazione della partizione di test di RAGTruth ha
individuato un numero considerevolmente maggiore di risposte allucinate rispetto alle etichette
originali~\cite{yu2026retromorphic}. L'insieme di test potrebbe quindi sottostimare le classi minoritarie. Per questo
motivo i risultati a tre classi su RAGTruth vanno sempre letti insieme a quelli su VitaminC, le cui etichette hanno un
accordo misurato. LLM-AggreFact unifica le etichette di undici dataset di qualità diversa, e alcune delle sue colonne sono
notoriamente rumorose.

\paragraph{Vantaggio nel dominio.} GroundOne è addestrato su RAGTruth, per cui i suoi risultati sull'insieme di test di
RAGTruth e sulla colonna RAGTruth di LLM-AggreFact lo favoriscono rispetto a sistemi addestrati altrove; per converso,
VitaminC rientra nella genealogia di addestramento di MiniCheck-RoBERTa. Ogni confronto è riportato con la relativa
avvertenza, e le affermazioni poggiano sull'andamento complessivo nei tre insiemi di test piuttosto che su uno solo.

\paragraph{Configurazione delle baseline.} I giudici generativi sono sensibili ai prompt. I loro prompt sono stati fissati
prima di esaminare i risultati di test, e il prompt del giudice a parità di output è stato scelto sui dati di sviluppo
unicamente in base alla validità del formato; l'interpretazione dell'output è stata permissiva a loro favore; la
decodifica non era vincolata da una grammatica, il che incide soprattutto sugli errori di formato del modello più piccolo.
Non è stato valutato un giudice addestrato sugli stessi dati, né un giudice più grande, per cui il confronto misura ciò
che offre un giudice della stessa dimensione senza addestramento. Il System One generalista è stato eseguito con pesi a 4
bit perché il modello a 8 bit non trovava posto accanto al sistema operativo, il che potrebbe sottostimarne
l'accuratezza. Il modello decisionale generalista d1-3B è stato eseguito zero-shot con le domande usate per tutte le
baseline System One, che potrebbero non essere la formulazione a esso più congeniale, e su un sottocampione di
LLM-AggreFact.

\paragraph{Contaminazione.} I modelli di base sono stati preaddestrati su ampi corpora web che potrebbero includere parti
degli insiemi di test pubblici. Poiché gli stessi modelli di base sono usati senza addestramento come baseline e
raggiungono un'accuratezza vicina al caso sul compito a tre classi, la contaminazione non può spiegare i guadagni
dell'addestramento, ma per i valori assoluti non può essere esclusa del tutto.
